\uuid{8rvI}
\titre{Intervalle}
\auteur{bodin ; hanani ; mzari}
\organisation{exo7}
\licence{CC BY-NC-SA 4.0}
\datecreate{2020-01-20}
\niveau{L1}
\module{Analyse}
\chapitre{Propriétés de R}
\sousChapitre{Intervalle, densité}
\difficulte{3}
\theme{Intervalle, densité, intervalle}
\exo7qcmid{lille-1.reels.027}

\contenu{
\begin{question}
\qclassification{120.04}
\qtags{motcle=intervalle}
Soit $I$ un intervalle ouvert de $\Rr$. Soient $x,y \in\Rr$ avec $x < y$.
Quelles sont les assertions vraies ?
\begin{answers}
    \good{Si $x,y\in I$, il existe $c\in I$ tel que $x < c < y$.}

    \good{Si $x,y\in I$, alors pour tout $c$ tel que $x < c < y$, on a $c \in I$.}

    \good{Si $x \notin I$ et $y\in I$, il existe $c\in I$ tel que $x < c < y$.}

    \bad{Si $x \notin I$ et $y\in I$, il existe $c\notin I$ tel que $x < c < y$.}
\end{answers}
\begin{explanations}
Si $x$ et $y$ sont deux éléments de l'intervalle $I$ alors toute valeur entre $x$ et $y$ est aussi dans l'intervalle.
\end{explanations}
\end{question}
}
